% =============================================================
% chapters/rozdil2.tex — РОЗДІЛ 2
% Розробка інформаційно-аналітичної системи моніторингу освітньої
% діяльності закладів професійної освіти.
%
% Структура (наповнення — Phase 2, T2.1..T2.10):
%   2.1 Ландшафт цифрових систем моніторингу освіти в Україні
%       2.1.1 ІСУО, ЄДЕБО, ДІСО, УЦОЯО, АІКОМ
%       2.1.2 EU4Skills, EU4VET, ETF цифрові ініціативи
%       2.1.3 Прогалина: обласний моніторинг ПТО + переміщені заклади
%   2.2 Вимоги до системи
%       2.2.1 Нормативно-правова основа (Закон №4574-IX, Закон «Про освіту»)
%       2.2.2 Функціональні вимоги (24 форми Spec v2.0 → 8 модулів)
%       2.2.3 Нефункціональні вимоги (a11y, audit, ABAC, perf)
%       2.2.4 Traceability-матриця: стаття закону → форма → модуль
%   2.3 Архітектура та проєктні рішення
%       2.3.1 Позиціонування у Design Science Research
%       2.3.2 Модульний моноліт: обґрунтування вибору
%       2.3.3 ABAC: декларативна мова політик доступу
%       2.3.4 Аудит з SHA-256 хеш-ланцюгом
%       2.3.5 Ідемпотентний прийом даних: compare-then-write
%       2.3.6 Multi-sheet Excel імпорт
%       2.3.7 Спеціальні рішення для переміщених закладів та ВПО-здобувачів
%   2.4 Емпірична евалюація
%       2.4.1 Audit chain integrity
%       2.4.2 ABAC correctness (матриця role × action × resource)
%       2.4.3 Excel round-trip fidelity
%       2.4.4 Latency / throughput
%       2.4.5 WCAG 2.1 AA compliance
%       2.4.6 Експертне оцінювання (n>=3, SUS + домен)
%       2.4.7 Сценарний walkthrough та кейс переміщеного закладу
%   2.5 Рекомендації та дорожня карта впровадження
%   Висновки до розділу 2
% =============================================================

\section{\MakeUppercase{Розробка інформаційно-аналітичної системи моніторингу}}
\label{sec:rozdil2}

% TODO[T2.2]: Підрозділ 2.1 -- ландшафт.
\subsection{Ландшафт цифрових систем моніторингу освіти в Україні}

\subsubsection{ІСУО, ЄДЕБО, ДІСО, УЦОЯО, АІКОМ: функції та обмеження}
\subsubsection{EU4Skills, EU4VET та ETF цифрові ініціативи}
\subsubsection{Прогалина: обласний моніторинг ПТО з урахуванням переміщених закладів}

% TODO[T2.3]: Підрозділ 2.2 -- вимоги.
\subsection{Вимоги до системи}

\subsubsection{Нормативно-правова основа}
\subsubsection{Функціональні вимоги}
\subsubsection{Нефункціональні вимоги}
\subsubsection{Traceability-матриця: стаття закону → форма → модуль}

% TODO[T2.4]: Підрозділ 2.3 -- дизайн.
\subsection{Архітектура та проєктні рішення}

\subsubsection{Позиціонування у Design Science Research}
\subsubsection{Модульний моноліт: обґрунтування вибору}
\subsubsection{ABAC: декларативна мова політик доступу}
\subsubsection{Аудит з SHA-256 хеш-ланцюгом}
\subsubsection{Ідемпотентний прийом даних: compare-then-write}
\subsubsection{Multi-sheet Excel імпорт}
\subsubsection{Візуалізація схеми даних у~стилі JSON Crack}
% TODO[T2.4]: Для візуалізації реляційної моделі та JSON-AST політик ABAC
% використати підхід із https://jsoncrack.com/ — node-graph layout, де
% кожна таблиця/поліс — вузол, FK/посилання — спрямовані ребра.
% Реалізація через TikZ (inline) або через експорт JSON-схеми в SVG/PDF
% утилітою jsoncrack-cli або через перевикористання їхнього layout-алгоритму
% у Python (graphviz + custom node renderer). Inline-код у chapter, без
% окремого figures/.
\subsubsection{Дизайн-рішення для переміщених закладів та ВПО-здобувачів}

% TODO[T2.5--T2.8]: Підрозділ 2.4 -- евалюація.
\subsection{Емпірична евалюація системи}

\subsubsection{Audit chain integrity}
\subsubsection{ABAC correctness}
\subsubsection{Excel round-trip fidelity}
\subsubsection{Латентність та пропускна здатність}
\subsubsection{WCAG 2.1 AA compliance}
\subsubsection{Експертне оцінювання ($\geq 3$ експертів)}
\subsubsection{Сценарний walkthrough та кейс переміщеного закладу}

% TODO[T2.9]: Підрозділ 2.5 -- рекомендації.
\subsection{Рекомендації щодо впровадження та дорожня карта}

\section*{Висновки до розділу 2}
\addcontentsline{toc}{section}{Висновки до розділу 2}

% TODO[T2.10]: 4--6 висновків, що відповідають RQ2.1--RQ2.4.
